El ajuste obtiene $a_{\mathrm{ef}}=\num{4.6137}$ mm Al. El punto reservado de 125 kVp presenta un residuo de HVL1 de \num{1.346}\,\%. Estos residuos no se interpretan como incertidumbre clínica.
\begin{table}[htbp]\centering\small\caption{Ajuste condicionado de HVL1 (mm Al) y cambio del factor combinado a 10 cm de lado.}\resizebox{\textwidth}{!}{\begin{tabular}{rrrrr}\toprule
kVp nominal & HVL1 observada & HVL1 calculada & Residuo (\%) & Cambio $C$ (\%)\\\midrule
\num{50.0000} & \num{2.3200} & \num{2.3978} & \num{3.3517} & \num{2.6432}\\
\num{60.0000} & \num{2.9200} & \num{2.9202} & \num{0.0077} & \num{2.8138}\\
\num{80.0000} & \num{3.9275} & \num{3.8531} & \num{-1.8949} & \num{2.7555}\\
\num{100.0000} & \num{4.8700} & \num{4.7952} & \num{-1.5358} & \num{2.4453}\\
\num{125.0000} & \num{5.8200} & \num{5.8983} & \num{1.3462} & \num{1.9823}\\
\bottomrule\end{tabular}}\end{table}
\begin{figure}[p]\centering\begin{tikzpicture}
\begin{axis}[dermaplot,name=hvl,height=6.3cm,title={A · Observaciones y modelo condicionado},xlabel={Tensión medida media por grupo (kVp)},ylabel={HVL1 (mm Al)},xmin=45,xmax=130,ymin=1.8,ymax=6.5,legend style={at={(.02,.98)},anchor=north west}]
\addplot[azul,only marks,mark=*,mark size=2.5pt]coordinates{(49.8,2.32)(60.4,2.92)(80.275,3.9275)(100,4.87)};\addlegendentry{Grupos usados en el ajuste}
\addplot[ocre,only marks,mark=diamond*,mark size=3pt]coordinates{(124.6,5.82)};\addlegendentry{Grupo reservado}
\addplot[turquesa,dashed,mark=square*,mark size=1.7pt]coordinates{(49.8,2.39776)(60.4,2.92023)(80.275,3.85308)(100,4.79521)(124.6,5.89835)};\addlegendentry{Modelo condicionado}
\end{axis}
\begin{axis}[dermaplot,at={(hvl.below south west)},anchor=north west,yshift=-1.1cm,height=4.3cm,xlabel={Tensión medida media por grupo (kVp)},ylabel={Residuo (\%)},xmin=45,xmax=130,ymin=-3,ymax=4.5,xtick={50,60,80,100,125}]
\addplot[gray,dashed]coordinates{(45,0)(130,0)};
\addplot[azul,only marks,mark=*,mark size=2.5pt]coordinates{(49.8,3.3517)(60.4,.0077)(80.275,-1.8949)(100,-1.5358)};
\addplot[ocre,only marks,mark=diamond*,mark size=3pt]coordinates{(124.6,1.3462)};
\end{axis}
\end{tikzpicture}

\caption{Compatibilidad del modelo condicionado con HVL1: valores y residuos $100(H_{\rm calc}/H_{\rm obs}-1)$. El grupo de 125 kVp se reserva; los cuatro registros de 80 kVp se agrupan. Sin barras de error porque no se dispone de incertidumbres. Las líneas solo guían la lectura entre calidades.}\label{fig:ajuste_hvl}\end{figure}
